\chapter{Software and reproducibility}
\label{ch:software}
\chapterpromise{Numerical responsibilities of the implementation, independent correctness tests,
model and dataset identifiers, and the requirements for regenerating the reported quantities.}

\section{Implementation architecture}
The package separates observation-family calculations from partition inference. A segment routine
computes the log integrated likelihood, any required posterior moments and, for nonconjugate
models, a clearly identified numerical approximation. The partition routine combines those
quantities with the normalized prior over admissible boundaries. Two recursions serve different
purposes: sum-product for integrated evidence and posterior marginals, and max-sum for the joint
MAP partition. Replicate and group models construct pooled or responsibility-weighted segment
scores before invoking the same recursion. Prediction must dispatch according to the fitted
observation family, rather than falling back to an unrelated sampling law.



\section{Independent checks of the exact recursions}
For short sequences, enumerate the full admissible partition set. At each fixed segment count,
compare the sum of prior-weighted block likelihood products against the forward and backward
recursions. Verify the partition normalizer separately using unit block scores. Check that the
ordered boundary and boundary-event marginals sum to one and $k-1$, respectively. Compare
posterior moments and global MAP decisions against direct enumeration, including tied maxima,
zero-probability blocks, negative Gaussian segment means and minimum-length constraints.
No self-consistency identity replaces an independent numerical oracle.

\section{Observation-model semantics}
Known Gaussian precisions, Poisson exposures and Binomial trial counts enter different likelihoods.
Their sufficient-statistic and base-measure terms must be derived from the sampling model. In
particular, raising a normalized observation density to an arbitrary positive power is not in
general equivalent to changing its precision or exposure. If likelihood powers are used,
posterior calculations must be described as generalized-Bayes inference rather than as ordinary
sampling-model evidence. Explicit masks distinguish missing observations from observations with
invalid or zero-valued descriptors.

\section{Numerical approximation diagnostics}
A nonconjugate segment routine must return the approximation method, the value of its objective
or evidence surrogate, and a convergence indicator. For each admissible test segment, compare
log scores and posterior moments against independent adaptive quadrature when sufficiently
low-dimensional. For a uniform log-score bound $\varepsilon$, global posterior-odds perturbation
bounds are conditional on that bound holding on \emph{every reachable segment}. Small empirical
error on selected segments alone does not establish such a hypothesis. Labels such as
Jaakkola--Jordan and P\'olya--Gamma variational Bayes must not be reported as separate
implementations when they share one computational path.

\section{Data, output and verification records}
A reproducible experiment identifies the original data source and observation model, the ordered
coordinate axis, data and descriptor checksums, train/validation/test partitions, model
hyperparameters, algorithm configuration, random seed and implementation version. Store both
successful estimates and failures, including the statistical denominator for reported
proportions and the numerical conditions behind timeouts. Synthetic fallback is disallowed in
analyses declared to use an external empirical dataset.

\section{Software availability}
The \BayesBreak package is published under the MIT license and available through
three channels:
\begin{itemize}
  \item \textbf{Source and issues:} \href{\ProjectGitHub}{\texttt{github.com/osolari/bayesbreak}}
  \item \textbf{Package index:} \href{\ProjectPyPI}{\texttt{pip install bayesbreak}}
  \item \textbf{Archived release:} \href{\ProjectDOIURL}{doi:\ProjectDOI} (Zenodo)
\end{itemize}
The preprint describing the method is available at
\href{\ProjectArXiv}{\texttt{arxiv.org/abs/2603.14681}}.

\section{Execution and publication controls}
The controlled experiments E00--E10 in Appendix~\ref{app:experiments} are defined with their
replication units, estimands, comparison methods and acceptance tests. Completed figures must
be regenerated only from verified run records. Figures produced under an observation model later
corrected in the implementation require refitting, not a change in caption or a numerical patch.
The single-column paper and this book shall use the same accepted results. A successful LaTeX
compilation is a typesetting check, not a scientific verification of the reported numerical data.
